\section{Gauge transfer and two first-prime source tests}\label{app:gauge}
\subsection{Winding a holonomy is different from transfer evolution}
Two-dimensional heat-kernel Yang--Mills supplies a rigorous control
with explicit Wilson functionals \cite{YM2}. Normalize the positive
$SU(2)$ Casimir so that the spin-$j$ representation has value $j(j+1)$.
For dimensionless area-coupling parameter $s>0$,
\begin{equation}\label{eq:heat-character}
 \mathbb E_s[\chi_j(U)]=(2j+1)e^{-sj(j+1)}.
\end{equation}
The normalized fundamental loop wound $m$ times is
$w_m(U)=\tfrac12\Tr_{\C^2}(U^m)=\cos(m\vartheta)$, for eigenvalues
$e^{\pm i\vartheta}$. The character identity
$2\cos(m\vartheta)=\chi_{m/2}-\chi_{(m-2)/2}$ for $m\geq2$ gives
\begin{equation}\label{eq:windings}
 \mathbb E_s[w_1]=e^{-3s/4},\qquad
 \mathbb E_s[w_m]=\frac12\left[(m+1)e^{-sm(m+2)/4}
                           -(m-1)e^{-sm(m-2)/4}\right].
\end{equation}
If the first return is set to $r\in(0,1)$, the second is
$(3r^{8/3}-1)/2$, strictly less than $r^2$. Indeed, for $y=r^{2/3}$,
\begin{equation}\label{eq:winding-failure}
 3y^4-2y^3-1=(y-1)(3y^3+y^2+y+1)<0.
\end{equation}
No overall primitive-length coefficient repairs this relative failure.

The same literal winding rule fails for every central probability
convolution semigroup with proposed moments
$\mathbb E_t[w_1]=r_t=e^{-\lambda t}$ and
$\mathbb E_t[w_2]=r_t^2$, $\lambda>0$. The normalized spin-one
character coefficient would be $c_1(t)=(1+2r_t^2)/3$. Convolution
requires multiplicativity, but
\begin{equation}\label{eq:semigroup-failure}
 c_1(t+u)-c_1(t)c_1(u)
 =\frac29(1-r_t^2)(1-r_u^2)>0\qquad(t,u>0).
\end{equation}

In contrast, let $H=\tfrac23C_2$ on the character Hilbert space with
Haar inner product, and choose the unit vector $\eta=\chi_{1/2}$.
Then $H\eta=\tfrac12\eta$ and
\begin{equation}\label{eq:transfer-return}
 \ip\eta{e^{-m(\log p)H}\eta}=p^{-m/2}.
\end{equation}
This propagates a fixed representation for a longer interval; it does
not wind the same holonomy. It passes the scalar repetition test.

Even so, the literal affine Euclidean source
$\Phi_L(f)=\int_{I_L}f(x)e^{-(b+x)H}\eta\dd x$, $b>L/2$,
has a Hankel kernel $C(2b+x+y)$, where $C(t)=\ip\eta{e^{-tH}\eta}$.
If $H\geq0$ and this kernel is invariant under small simultaneous
translations of $x,y$, then $C'$ vanishes on an interval. Its spectral
measure satisfies
$-C'(t)=\int\lambda e^{-t\lambda}\dd\nu_\eta(\lambda)$, so
$\eta$ must lie in the zero-energy subspace. The resulting kernel is
constant. This excludes that affine-time source, not other spatial,
normalized or extended gauge preparations.

\subsection{A positive disk system for masses and returns}
Let $\Lambda$ be the Dirichlet-to-Neumann operator on the unit disk.
On the reflection-even boundary subspace, with normalized measure
$\dd\vartheta/(2\pi)$, use the basis
$e_0=1$, $e_n=\sqrt2\cos(n\vartheta)$ for $n\geq1$.
Harmonic extension gives $\Lambda e_n=ne_n$. Thus
\begin{equation}\label{eq:disk-A}
 A=2\Lambda+\tfrac12\id,
 \qquad Ae_n=a_ne_n,
\end{equation}
has the gamma spectrum and positive quadratic form
\begin{equation}\label{eq:disk-energy}
 \ip h{Ah}=\frac2{2\pi}\int_{\mathbb D}|\nabla u_h|^2\dd^2z
       +\frac{1/2}{2\pi}\int_{S^1}|h|^2\dd\vartheta.
\end{equation}
Here $u_h$ is the harmonic extension; the form domain is the even
$H^{1/2}$ trace space. The choice of half-unit shift and one mode
per nonnegative integer is explicit model data, not forced by an
arithmetic symmetry.

The transfer is a shifted Poisson semigroup. On the full circle it
also has the positive subordination formula
\begin{equation}\label{eq:subordination}
 e^{-tA}=e^{-t/2}\frac t{\sqrt\pi}
 \int_0^\infty s^{-3/2}e^{-t^2/s}e^{-sD^2}\dd s,
 \qquad D=-i\partial_\vartheta,
\end{equation}
which restricts to the even subspace. The scalar Laplace integral
proves the formula on each Fourier mode. This is a change of transfer
generator, not an equivalence with the original Yang--Mills dynamics.

Let $P_0=|e_0\rangle\langle e_0|$ and $T_p=e^{-(\log p)A}$.
Then
\begin{align}
 b(\tau^2)&=2\Tr\left[A^{-1}\frac{\tau^2}{A^2+\tau^2}\right],
 \notag\\
 \Tr(P_0T_p^m)&=p^{-m/2},\qquad
 T_pT_\ell=T_{p\ell},\qquad
 \Tr(T_p^m)=\frac{p^{-m/2}}{1-p^{-2m}}.
 \label{eq:disk-transfer}
\end{align}
The ground projection is therefore essential to the proposed pure
geometric return. This model does not derive primality, the primitive
coefficient $\log p$, or the absence of mixed histories.

\subsection{An orthogonal first-prime branch has an unavoidable cusp}
Use the source \eqref{eq:gamma-source} and set $a=1/2$,
$d=\log2$, $r=e^{-ad}=2^{-1/2}$. With $H_x=-\partial_x^2$,
consider the full positive preparation
\begin{equation}\label{eq:branch-source}
 F\longmapsto\left(
 \{\mathcal A_nF\}_{n\geq1},
 \sqrt{1-\theta}\mathcal A_0F,
 \sqrt\theta\mathcal A_0(\id-rU_d)F\right),
 \qquad 0\leq\theta\leq1.
\end{equation}
Its Gram operator is, in form notation,
$K+\theta(2/a)H_x(a^2+H_x)^{-1}
[r^2\id-r(U_d+U_d^*)]$.
Matching the prime atom forces $\theta=ad/2=d/4$. The resulting
complete pairing has operator
\begin{equation}\label{eq:branch-gram}
 B=K+dr^2\id-dr(U_d+U_d^*)+R_{a,d},
 \qquad
 R_{a,d}=-da^2[r^2\id-r(U_d+U_d^*)](a^2+H_x)^{-1}.
\end{equation}
The resolvent kernel $e^{-a|s|}/(2a)$ gives
\begin{equation}\label{eq:cusp-kernel}
 R_{a,d}(s)=\frac{ad}{2}
 \begin{cases}
 r^2e^{a|s|},&|s|\leq d,\\
 e^{-a|s|},&|s|\geq d.
 \end{cases}
 \qquad
 R_{a,d}'(d+)-R_{a,d}'(d-)=-da^2r\ne0.
\end{equation}
On $d<L<\log3$, after the common prime atom is removed,
the candidate-minus-target difference is
$(dr^2-w_0)\id+E_L^*R_{a,d}E_L-P_L$. The pole kernel is smooth
near $x-y=d$ and the contact has no support there. Neither can
cancel the cusp, so the full pairing cannot match on that interval.
This is an exact kernel obstruction, requiring no residual bound.

For a finite positive sum of analogous orthogonal branches with mass
$a_n$, positive weight $\theta_n$ and attenuation $r_n$, the derivative
jump at $d$ is $-2\theta_na_nr_n$ for each active branch. The signs
agree and cannot cancel. General coherent changes of source are
outside this assertion.

\subsection{A coherent repair cancels the cusp and still fails the contact}
In two common outputs, use the lowest gamma source and a shift source
\begin{equation}\label{eq:coherent-sources}
 \Gamma_0F=\begin{pmatrix}
  2H_x(a^2+H_x)^{-1}F\\ \partial_x(a^2+H_x)^{-1}F
 \end{pmatrix},\quad
 Z_{u,v}F=\begin{pmatrix}
 u(U_d-U_d^*)F\\vF+u(U_d+U_d^*)F
 \end{pmatrix},\quad a=\tfrac12,
\end{equation}
where $u,v$ are real. All operators commute on the line. The first
components are respectively self-adjoint and skew-adjoint, and the
second components respectively skew-adjoint and self-adjoint. Hence
\begin{equation}\label{eq:coherent-cross}
 \Gamma_0^*Z+Z^*\Gamma_0=0,
 \qquad Z^*Z=(v^2+4u^2)\id+2uv(U_d+U_d^*).
\end{equation}
Keeping the other gamma modes and setting $2uv=-dr$ gives
\begin{equation}\label{eq:coherent-result}
 B_{\rm coh}=K+C\id-dr(U_d+U_d^*),\qquad
 C=v^2+4u^2\geq2dr.
\end{equation}
The inequality follows from $C-2dr=(v+2u)^2$. There is no displaced
continuous term, but on the first-prime interval
\begin{equation}\label{eq:coherent-failure}
 B_{{\rm coh},L}-Q_L=(C-w_0)\id-P_L\ne0.
\end{equation}
Indeed $C-w_0>0$ while $P_L$ has rank two on an infinite-dimensional
input space. This coherent representation is still the sum of independent
gamma and jump Gram forms, because its Hermitian interference vanishes.
It supplies a second exact stopping criterion for the tested source.
